%% kapa: exam=1 round=34 year=2023 period=1 form=A subject=civil_law subject_ko=민법 seq=33 number=33
%% source: Q-Net 감정평가사 기출문제 (cst00309) artlSeq=5211789 file=raw/1008/5211789/2023년도 제34회 감정평가사 1차시험 1교시.pdf
%% conversion: kordoc 4.24.51 -> tools/gen1.py (auto); review=proofread-sonnet; flags=none
\begin{question}{33}
지상권에 관한 설명으로 옳은 것은? (다툼이 있으면 판례에 따름)
\begin{choices}
  \item 건물의 소유를 목적으로 하는 지상권의 양도는 토지소유자의 동의를 요한다.
  \item 지료합의가 없는 지상권 설정계약은 무효이다.
  \item 수목의 소유를 목적으로 하는 지상권의 최단존속기간은 10년이다.
  \item 지상권이 설정된 토지의 소유자는 그 지상권자의 승낙 없이 그 토지 위에 구분지상권을 설정할 수 있다.
  \item 「장사 등에 관한 법률」 시행 이전에 설치된 분묘에 관한 분묘기지권의 시효취득은 법적 규범으로 유지되고 있다.
\end{choices}
\end{question}
\answer{5}
